% =========================================================================
% Capítulo 4: Cuerpo rígido
% Tema 01: Cuerpo rígido
% Problema 06: Momentos de inercia de figuras 2D y tensor de inercia 3D
% =========================================================================
\subsubsection{Momentos de inercia de figuras 2D y tensor de inercia 3D}

El cálculo de momentos y productos de inercia es la base del análisis dinámico de cuerpos
rígidos. En este ejercicio se pide obtener, por integración directa, las propiedades de
inercia de varios cuerpos homogéneos.

\begin{enumerate}
    \item \textbf{Figuras 2D} (densidad constante, masa $M$). Determine el momento de inercia
          respecto al eje perpendicular al plano que pasa por el centro de masa para:
          \begin{itemize}
              \item una circunferencia (aro delgado) de radio $R$;
              \item un disco (círculo lleno) de radio $R$;
              \item una barra delgada de largo $L$.
          \end{itemize}
    \item \textbf{Tensor de inercia 3D} (densidad constante). Determine el tensor de inercia
          respecto al centro de masa para:
          \begin{itemize}
              \item un cascarón esférico de radio $R$ y masa $M$;
              \item un paralelepípedo de dimensiones $L\times H\times W$ y masa $M$.
          \end{itemize}
\end{enumerate}

Recuerde que el tensor de inercia se define como
\[
\mathbf{I} =
\begin{bmatrix}
\displaystyle\int_V \rho(y^2 + z^2)\, dV & \displaystyle -\int_V \rho\, x y\, dV & \displaystyle -\int_V \rho\, x z\, dV \\[4pt]
\displaystyle -\int_V \rho\, x y\, dV & \displaystyle\int_V \rho(x^2 + z^2)\, dV & \displaystyle -\int_V \rho\, y z\, dV \\[4pt]
\displaystyle -\int_V \rho\, x z\, dV & \displaystyle -\int_V \rho\, y z\, dV & \displaystyle\int_V \rho(x^2 + y^2)\, dV
\end{bmatrix}.
\]

% --- SOLUCIÓN ---------------------------------------------------------------
\begin{solucion}{Momentos de inercia de figuras 2D y tensor de inercia 3D}
\subsubsection*{Pregunta 1: Figuras 2D}

\textbf{Circunferencia (aro delgado) de radio $R$.} Toda la masa está a distancia $R$ del
eje central perpendicular al plano:
$$I = \int R^2\,dm = R^2\!\int dm = \boxed{MR^2}.$$

\textbf{Disco de radio $R$.} Con densidad superficial $\sigma=M/(\pi R^2)$ y $dm=\sigma\,2\pi\rho\,d\rho$:
$$I = \int_0^R \rho^2\,\sigma\,2\pi\rho\,d\rho = \sigma\,2\pi\frac{R^4}{4}
= \frac{M}{\pi R^2}\,\frac{\pi R^4}{2} = \boxed{\tfrac12 MR^2}.$$

\textbf{Barra delgada de largo $L$.} Con densidad lineal $\lambda=M/L$, respecto a un eje
perpendicular que pasa por su centro:
$$I = \int_{-L/2}^{L/2} x^2\,\lambda\,dx = \frac{M}{L}\,\frac{2}{3}\left(\frac{L}{2}\right)^3
= \boxed{\tfrac{1}{12}ML^2}.$$
(Por el teorema de Steiner, respecto a un extremo, $I=\tfrac13 ML^2$.)

\subsubsection*{Pregunta 2: Tensor de inercia 3D}

\textbf{Cascarón esférico de radio $R$.} Por simetría todos los productos de inercia se
anulan y los tres momentos principales son iguales. Como $x^2+y^2+z^2=R^2$ sobre el cascarón,
$$I_{xx}+I_{yy}+I_{zz} = \int 2(x^2+y^2+z^2)\,dm = 2R^2 M \;\Rightarrow\; I_{xx}=I_{yy}=I_{zz}=\tfrac23 MR^2,$$
de modo que
$$\mathbf{I} = \frac{2}{3}MR^2
\begin{bmatrix} 1 & 0 & 0 \\ 0 & 1 & 0 \\ 0 & 0 & 1 \end{bmatrix}.$$

\textbf{Paralelepípedo $L\times H\times W$.} Con los ejes centrados y paralelos a las aristas
($L$ en $x$, $H$ en $y$, $W$ en $z$), los productos de inercia se anulan por simetría y
$$\int_V x^2\,dV = \frac{L^2}{12}M,\quad \int_V y^2\,dV = \frac{H^2}{12}M,\quad \int_V z^2\,dV = \frac{W^2}{12}M,$$
por lo que
$$\mathbf{I} = \frac{M}{12}
\begin{bmatrix} H^2+W^2 & 0 & 0 \\ 0 & L^2+W^2 & 0 \\ 0 & 0 & L^2+H^2 \end{bmatrix}.$$
\end{solucion}
